\begin{description}
\item[(a)] The longitude of Gaurishankar is
\[\lambda_g=\frac{5.75}{12}\times180^\circ=86.25^\circ\mathrm E.\]
\hfill(3 points)
The official calculation gives
\[\lambda_k=\lambda_g-\frac{105\cos\phi}{6378}\frac{180^\circ}{\pi}=85.415^\circ\mathrm E.\]
\hfill(4 points)
% sic: The official key multiplies by cos(phi). A longitudinal displacement along the latitude circle instead requires division by cos(phi).
\item[(b)] The official time difference is
\[t_{\mathrm{diff}}=\frac{105\cos\phi}{6378}\frac{12\times3600}{\pi}\,\mathrm s=200.4\,\mathrm s=3\,\mathrm{min}\ 20.4\,\mathrm s.\]
\hfill(3 points)
Then $t_{\mathrm{Kathmandu}}=8{:}35\,\mathrm{pm}-t_{\mathrm{diff}}=8{:}31{:}40\,\mathrm{pm}$. Answers of 8:32 pm or 8:31:40 pm are accepted. \hfill(1 point)
\item[(c)] Ravi has the same local time as Kathmandu. On his watch, however, his time zone is UTC $+5{:}30$. Thus the difference is $15-t_{\mathrm{diff}}$ minutes, or 11 minutes 40 (39.6) seconds. \hfill(2 points)
As per his watch, therefore, the occultation will occur at 8:16:40 pm, or 8:17 pm. \hfill(2 points)
% sic: The source calls the meridian crossing an occultation and gives 8:16:40 pm, inconsistent with the stated 15-minute difference between the two watches. Values and wording are retained from the official solution.
\end{description}
